% Generated draft: requires mathematical and source review.
\documentclass[11pt]{article}
\usepackage[margin=25mm]{geometry}
\usepackage{fontspec}
\setmainfont{Latin Modern Roman}
\usepackage{amsmath,amssymb,mathtools,unicode-math}
\setmathfont{Latin Modern Math}
\usepackage{enumitem,needspace,xurl,hyperref}
\setlength{\parindent}{0pt}
\setlength{\parskip}{5pt}
\setlength{\emergencystretch}{3em}
\newcommand{\N}{\mathbb{N}}
\newcommand{\Z}{\mathbb{Z}}
\newcommand{\Q}{\mathbb{Q}}
\newcommand{\R}{\mathbb{R}}
\newcommand{\C}{\mathbb{C}}
\newcommand{\F}{\mathbb{F}}
\newcommand{\A}{\mathbb{A}}
\newcommand{\PP}{\mathbb{P}}
\newcommand{\E}{\mathbb{E}}
\newcommand{\eps}{\varepsilon}
\newcommand{\dd}{\,\mathrm{d}}
\begin{document}
\section*{UnsolvedMath ID: 20002896}
\textbf{Source problem code:} catalogue entry
\subsection{Definitions and mathematical statement}
An \(r\)-ribbon is a connected skew Young diagram of \(r\) cells containing no \(2\times2\) square. Fix a ribbon length \(r\ge1\), a finite alphabet \(\{1,\ldots,N\}\), a ribbon-tileable skew shape, and a convention for semistandard ribbon tableaux, equivalently the usual chains of shapes adding horizontal \(r\)-ribbon strips of each label. Let \(\mathcal T\) be this fixed family. Its content weight is \(\operatorname{wt}(T)=(m_1,\ldots,m_N)\), where \(m_i\) counts ribbons labeled \(i\). In the half-integer spin convention,
\[
\operatorname{spin}(T)=\frac12\sum_{R\text{ a ribbon of }T}(\operatorname{height}(R)-1),
\]
where height is the number of occupied rows; doubling this statistic is another common normalization.

The AIM problem asks for a crystal structure on ribbon tableaux compatible with spin. A precise type-\(A\) interpretation is to construct operators
\[
e_i,f_i:\mathcal T\longrightarrow\mathcal T\cup\{\varnothing\},\qquad 1\le i<N,
\]
with content as weight, the usual inverse and weight-change rules, and a regular normal finite type-\(A_{N-1}\) crystal structure: its connected components are isomorphic to ordinary semistandard tableau crystals \(B_N(\nu)\). Require
\[
f_iT\ne\varnothing\Longrightarrow\operatorname{spin}(f_iT)=\operatorname{spin}(T),
\]
and the analogous condition for \(e_i\). Thus spin is constant on each crystal component. The source states that the type-\(A\) domino case \(r=2\) has been solved; it does not give its reference or define the precise intended compatibility convention.

The supplied research aid reports a positive special case for every \(r\) and finite alphabet: a fixed semistandard ribbon family whose \(r\)-quotient has at most two nonempty components, each consisting of a single cell, admits an explicit content-weight regular type-\(A\) crystal preserving spin, cospin and diagonal inversions on each component. Cospin may be normalized as the maximum spin for the fixed shape minus spin. With two cells, after choosing the attacking order of their labels, the weak-order sector is \(B_N((2))\), and the inversion sector is \(B_N((1,1))\); for \(N=1\) the latter sector is empty. The empty and one-cell quotients give the corresponding trivial and one-box crystals.

The report emphasizes that every pair of distinct unicellular quotient components attacks in its quotient model. Consequently the realizable matching case has at most one edge, exactly the stated at-most-two-cell setting; no independent multi-edge factorization is claimed. It also reports a transport construction from a statistic-preserving strong dual-equivalence component together with its standardization fibers to an ordinary-tableau crystal. The full hypotheses of that more general transport statement are not supplied here. The unrestricted ribbon-tableau crystal problem is not solved by these special cases.
\subsection{Short English statement}
Construct content-weight crystals on ribbon tableaux whose crystal moves preserve spin. The source records the domino case as solved, and the supplied aid treats quotients with at most two singleton components; it does not solve the unrestricted problem or claim a multiple-edge factorization.
\subsection{Sources}
\begin{itemize}
\item AIM workshop \emph{Generalized Kostka polynomials}, Problem (1), \url{https://aimath.org/WWN/kostka/prob2.pdf}. The source states that the type-\(A\) domino case is solved but supplies no separate bibliography.
\item Local UnsolvedMath registry, permanent ID 20002896, exact code \texttt{catalogue entry}; canonical location \texttt{aim-special-functions-notes.json}, notes[0], dated 14 August 2026; statement labeled exact.
\item Research attempt 1: partial result, problem partially solved at the run, rigor label proved, novelty confidence low; registry status partially solved. The two-cell and transport constructions are supplied machine-generated research aids, not peer reviewed, requiring independent verification before citation.
\end{itemize}
\end{document}
